\section*{अध्याय \ref{ch:b1:geological-maps} --- भूवैज्ञानिक मानचित्र}

\begin{solution}{exo:b1:geological-maps:1}
\omterm{def:b1:geological-maps:attitude}{स्ट्राइक}: समतल की क्षैतिज रेखाओं का दिगंश; \omterm{def:b1:geological-maps:attitude}{नति}: समतल और क्षैतिज के बीच का कोण, जो \omterm{def:b1:geological-maps:attitude}{स्ट्राइक} के लंबवत मापा
जाता है; \omterm{def:b1:geological-maps:attitude}{नति-दिशा}: सबसे खड़ी अधोमुखी रेखा का दिगंश। $135/40$ के लिए, दाएँ-हाथ के नियम से \omterm{def:b1:geological-maps:attitude}{स्ट्राइक}
$135 - 90 = \ang{045}$ है। चिह्न: एक \omterm{def:b1:geological-maps:attitude}{स्ट्राइक} रेखा जो उत्तर से \ang{45} पूर्व की ओर अभिमुख है, जिसके दक्षिण-पूर्वी
भाग पर (\ang{135} की ओर) एक निशान है, जिस पर 40 अंकित है।
\end{solution}

\begin{solution}{exo:b1:geological-maps:2}
$z = 500 - \tan\ang{30}\,(x\sin\ang{90} + y\cos\ang{90}) = 500 - 0.577\,x$।
$(1000, 0)$ पर: $500 - 577 = \qty{-77}{m}$। $(0, 1000)$ पर: \qty{500}{m}, क्योंकि समतल अपनी उत्तर-दक्षिण \omterm{def:b1:geological-maps:attitude}{स्ट्राइक}
के अनुदिश क्षैतिज है।
\end{solution}

\begin{solution}{exo:b1:geological-maps:3}
\qty{350}{m} की समोच्च रेखा का अनुसरण करने वाला अनुरेख एक क्षैतिज \omterm{def:b1:geological-maps:geological-map}{संपर्क} है। पहाड़ियों और घाटियों के आरपार सीधी
रेखा में चलने वाला अनुरेख एक ऊर्ध्वाधर \omterm{def:b1:geological-maps:geological-map}{संपर्क} का है।
\end{solution}

\begin{solution}{exo:b1:geological-maps:4}
नदी के किनारे \omterm{def:b1:geological-maps:fold}{अपनति} के क्रोड में शिस्ट और ग्रेनाइट (बलुआ पत्थर के नीचे) हैं; \omterm{def:b1:geological-maps:fold}{अभिनति} के क्रोड में खड़िया है।
सबसे पुरानी इकाइयाँ अनुक्रम में सबसे नीचे रहती हैं: जहाँ परतें ऊपर की ओर मेहराब बनाती हैं और अपरदन उनमें
सबसे गहरा काटता है, वहाँ सबसे पुरानी इकाइयाँ दिखती हैं; जहाँ वे झुकती हैं, वहाँ सबसे नई इकाइयाँ बीच में
संरक्षित रहती हैं।
\end{solution}

\begin{solution}{exo:b1:geological-maps:5}
पहले दो बिंदु समान उन्नयन पर हैं: रेखा $y = 0$ एक \omterm{def:b1:geological-maps:attitude}{स्ट्राइक} रेखा है, अतः \omterm{def:b1:geological-maps:attitude}{स्ट्राइक} पूर्व-पश्चिम है। \omterm{def:b1:geological-maps:geological-map}{संपर्क} उत्तर की
ओर \qty{500}{m} पर \qty{100}{m} नीचे है: वह उत्तर की ओर झुकता है, $\tan\delta = 100/500 = 0.2$,
$\delta \approx \ang{11.3}$; \omterm{def:b1:geological-maps:attitude}{नति-दिशा} \ang{000}, दाएँ-हाथ के नियम से \omterm{def:b1:geological-maps:attitude}{स्ट्राइक} \ang{270}। प्रमेय से:
$\overrightarrow{P_1P_2} = (1000, 0, 0)$,
$\overrightarrow{P_1P_3} = (0, 500, -100)$, $\vec n = (0, 100\,000,
500\,000)$, $\cos\delta = 500\,000/509\,902 \approx 0.981$, $\delta \approx
\ang{11.3}$, और $(n_x, n_y) = (0, 100\,000)$ उत्तर की ओर इंगित करता है: वही परिणाम।
\end{solution}

\begin{solution}{exo:b1:geological-maps:6}
दिशा प्रवाह में उत्तर है, $s = 0.03$ (\ang{1.72})। (क) $a = 0 < s$: V प्रवाह के विपरीत (दक्षिण), समोच्चों के
अनुदिश। (ख) $a = -\tan\ang{20} < 0 < s$: V प्रवाह के विपरीत। (ग) $a = \tan\ang{1} \approx 0.017 < s$: V प्रवाह
के विपरीत। (घ) $a = \tan\ang{20}
\approx 0.364 > s$: V प्रवाह की ओर (उत्तर)।
\end{solution}

\begin{solution}{exo:b1:geological-maps:7}
$\xi_{\text{तल}} - \xi_{\text{किनारा}} = h/(a - s) = 80/(0.364 - 0.03)
\approx \qty{240}{m}$: किनारे पर अनुरेख शीर्ष के लगभग \qty{240}{m} प्रवाह के विपरीत है।
\end{solution}

\begin{solution}{exo:b1:geological-maps:8}
\omterm{def:b1:geological-maps:fold}{अपनति} का बंद होना नमन की दिशा में इंगित करता है: \omterm{def:b1:geological-maps:fold}{वलन-अक्ष} उत्तर की ओर नमन करता है। उत्तर की ओर नमन
करते \omterm{def:b1:geological-maps:fold}{वलन-अक्ष} के साथ एक \omterm{def:b1:geological-maps:fold}{अभिनति} के \omterm{def:b1:geological-maps:geological-map}{बहिर्भाग} दक्षिण की ओर बंद होंगे।
\end{solution}

\begin{solution}{exo:b1:geological-maps:9}
पूर्व-दक्षिण-पूर्व की ओर की सतह पर नई इकाइयाँ (सामने की पुरानी इकाइयों के विपरीत खड़िया और शेल) हैं: वह
खंड नीचे खिसका। भ्रंश-चिह्न के निशान उसी ओर हैं, अवनत ओर।
\end{solution}

\begin{solution}{exo:b1:geological-maps:10}
एक रेखा जो दोनों समतलों में है, दोनों अभिलंबों के लंबवत है, अतः $\vec n_1 \times \vec n_2$ के समांतर है। $090/30$
के लिए: $\vec n_1 = (0.577, 0,
1)$; $270/30$ के लिए: $\vec n_2 = (-0.577, 0, 1)$। तब $\vec n_1 \times \vec n_2
= (0\cdot 1 - 1\cdot 0,\ 1\cdot(-0.577) - 0.577\cdot 1,\ 0) = (0, -1.155,
0)$: एक क्षैतिज उत्तर-दक्षिण रेखा। पहला समतल पूर्व की ओर झुकता है, दूसरा पश्चिम की ओर: वे रेखा से दूर झुकते
हैं, जो क्षैतिज अक्ष वाले एक \omterm{def:b1:geological-maps:fold}{अपनति} का शिखर है।
\end{solution}

\begin{solution}{exo:b1:geological-maps:11}
ऊर्ध्वाधर डाइक का अनुरेख मूल-बिंदु से गुज़रने वाली \ang{045} दिशा की सीधी रेखा है, किसी भी उन्नयन पर; पहाड़ी पर वह
उसी रेखा पर कहीं है। नत डाइक समतल $z = -\tan\ang{60}\,d$ है, जहाँ $d$ मूल-बिंदु से गुज़रने वाली \omterm{def:b1:geological-maps:attitude}{स्ट्राइक} रेखा से
क्षैतिज दूरी है, जो दक्षिण-पूर्व की ओर धनात्मक गिनी जाती है। $z = \qty{300}{m}$ पर: $d = -300/\tan\ang{60}
\approx \qty{-173}{m}$: अनुरेख मूल-बिंदु से गुज़रने वाली \omterm{def:b1:geological-maps:attitude}{स्ट्राइक} रेखा से \qty{173}{m} उत्तर-पश्चिम में है, \omterm{def:b1:geological-maps:attitude}{नति} के
ऊपर-ढलान वाली ओर।
\end{solution}

\begin{solution}{exo:b1:geological-maps:12}
\[
  \frac{\partial z}{\partial x} \approx \frac{405 - 420}{100} = -0.15, \qquad
  \frac{\partial z}{\partial y} \approx \frac{412 - 420}{100} = -0.08,
\]
अतः ढलान $\arctan\sqrt{0.15^2 + 0.08^2} = \arctan 0.17 \approx
\ang{9.6}$ है। ज़मीन $-\nabla z = (0.15, 0.08)$ के अनुदिश नीचे उतरती है, जिसका दिगंश
$\arctan(0.15/0.08) \approx \ang{62}$ है: घन पूर्व-उत्तर-पूर्व की ओर मुँह करता है।
\end{solution}

\begin{solution}{pb:b1:geological-maps:1}
\textbf{1.} शिस्ट, ग्रेनाइट, बलुआ पत्थर, चूना-पत्थर, शेल, खड़िया, जलोढ़।

\textbf{2.} पश्चिम से पूर्व: खड़िया और शेल, चूना-पत्थर, फिर बलुआ पत्थर, फिर $x = \qty{3.5}{km}$ के पास शिस्ट और
ग्रेनाइट, फिर बलुआ पत्थर, चूना-पत्थर और शेल फिर, फिर खड़िया; पूरे तल पर जलोढ़।

\textbf{3.} बलुआ पत्थर, जिसके नीचे नदी के किनारे शिस्ट और ग्रेनाइट हैं: क्रोड में सबसे पुरानी इकाइयाँ, एक \omterm{def:b1:geological-maps:fold}{अपनति}।

\textbf{4.} खड़िया, सबसे नई इकाई: एक \omterm{def:b1:geological-maps:fold}{अभिनति}।

\textbf{5.} बलुआ पत्थर का आधार \omterm{def:b1:geological-maps:fold}{अपनति} के अक्ष पर सबसे ऊँचा है, और केवल नदी के किनारे अपरदन ने उसके नीचे काटा है।

\textbf{6.} एक असंगति: बलुआ पत्थर दो भिन्न पुरानी इकाइयों, शिस्ट और ग्रेनाइट, पर टिकता है, और ग्रेनाइट की रूपरेखा
काट देता है।

\textbf{7.} $435 - 261 = 174$, $432 - 428 = 4$, $512 - 586 = \qty{-74}{m}$।

\textbf{8.} $\overrightarrow{E_1E_2} = (600, 1200, -170)$,
$\overrightarrow{E_1E_3} = (1100, 400, -248)$।

\textbf{9.} सदिश गुणनफल है
\[
  (-229\,600,\ -38\,200,\ -1\,080\,000);
\]
ऊपर की ओर अभिमुख, $\vec n_E = (229\,600,\ 38\,200,\ 1\,080\,000)$।

\textbf{10.} $\|\vec n_E\| \approx 1\,104\,800$; $\cos\delta \approx 0.9776$,
$\delta \approx \ang{12.2}$।

\textbf{11.} $(229\,600, 38\,200)$ का दिगंश: $\arctan(229\,600/38\,200)
\approx \ang{80.6}$; \omterm{def:b1:geological-maps:attitude}{अभिविन्यास} $081/12$, \omterm{def:b1:geological-maps:attitude}{स्ट्राइक} $\approx \ang{351}$।

\textbf{12.} हाँ: पूर्वी भुजा के चिह्न की \omterm{def:b1:geological-maps:attitude}{स्ट्राइक} रेखा लगभग उत्तर-दक्षिण है और नति-निशान पूर्व की ओर है, जिस पर 12 अंकित है।

\textbf{13.} $524 - 669 = -145$, $421 - 459 = -38$, $319 - 209 =
\qty{110}{m}$।

\textbf{14.} $\overrightarrow{W_1W_2} = (700, 1200, 107)$ और
$\overrightarrow{W_1W_3} = (1100, -600, 255)$। उनका सदिश गुणनफल, ऊपर की ओर अभिमुख, है
\[
  (370\,200,\ -60\,800,\ -1\,740\,000) \quad\longrightarrow\quad
  \vec n_W = (-370\,200,\ 60\,800,\ 1\,740\,000).
\]

\textbf{15.} $\cos\delta = 1\,740\,000/\|\vec n_W\| \approx 0.9776$, $\delta
\approx \ang{12.2}$; \omterm{def:b1:geological-maps:attitude}{नति-दिशा}: $(-370\,200, 60\,800)$ का दिगंश,
$\approx \ang{279.3}$। \omterm{def:b1:geological-maps:attitude}{अभिविन्यास} $279/12$, \omterm{def:b1:geological-maps:attitude}{स्ट्राइक} $\approx \ang{189}$: पश्चिमी भुजा पश्चिम-उत्तर-पश्चिम की ओर झुकती है,
पूर्वी भुजा पूर्व-उत्तर-पूर्व की ओर।

\textbf{16.} $\vec n_E \times \vec n_W = (804 \times 10^6,\ -799\,320 \times
10^6,\ 28\,101 \times 10^6)$ (पूर्णांकित)।

\textbf{17.} नीचे की दिशा $(-804, 799\,320, -28\,101) \times
10^6$ है: दिशा $\arctan(-804/799\,320) \approx \ang{0}$, यानी उत्तर; नमन
$\arctan\bigl(28\,101/\sqrt{804^2 + 799\,320^2}\bigr) \approx \ang{2.0}$। \omterm{def:b1:geological-maps:fold}{वलन-अक्ष} उत्तर की ओर \ang{2} नमन करता है।

\textbf{18.} $8000 \times \tan\ang{2.0} \approx \qty{280}{m}$ नीचे, लगभग \qty{-50}{m} के निकट। \omterm{def:b1:geological-maps:fold}{अपनति} का क्रोड नदी के
पास सबसे ऊँचा है और उत्तर की ओर धँसता है, इसलिए उसका \omterm{def:b1:geological-maps:geological-map}{बहिर्भाग} उत्तर की ओर बंद होता है, जैसा मानचित्र दिखाता है।

\textbf{19.} $\overrightarrow{F_1F_2} = (1627, 2500, -276)$ और
$\overrightarrow{F_1F_3} = (3124, 5500, 77)$ ऊपर का अभिलंब देते हैं
\[
  (1\,710\,500,\ -987\,503,\ 1\,138\,500), \qquad
  \delta = \arccos\frac{1\,138\,500}{2\,279\,700} \approx \ang{60},
\]
और \omterm{def:b1:geological-maps:attitude}{नति-दिशा} $\approx \ang{120}$। \omterm{def:b1:geological-maps:attitude}{अभिविन्यास} $120/60$, \omterm{def:b1:geological-maps:attitude}{स्ट्राइक} \ang{030}।

\textbf{20.} भ्रंश पूर्व-दक्षिण-पूर्व की ओर झुकता है, इसलिए वह ओर लटकती दीवार है, वह खंड जो झुके भ्रंश-तल के ऊपर है; वह नीचे खिसका (सतह पर नई इकाइयाँ): एक सामान्य भ्रंश।

\textbf{21.} भ्रंश-तल पूर्व-दक्षिण-पूर्व की ओर झुकता है, प्रवाह की दिशा में, और नदी के तल के उतार से कहीं अधिक
तीखी \omterm{def:b1:geological-maps:attitude}{नति} से: निचली ज़मीन तल को \omterm{def:b1:geological-maps:attitude}{नति} की दिशा में और नीचे मिलती है, इसलिए अनुरेख घाटी में पूर्व-दक्षिण-पूर्व की ओर
मुड़ता है --- भ्रंश-तल पर V-नियम।

\textbf{22.} ग्रेनाइट का अंतर्वेधन; बलुआ पत्थर से पहले का अपरदन; खड़िया का निक्षेपण; \omterm{def:b1:geological-maps:fold}{वलन}; भ्रंशन; जलोढ़ का निक्षेपण।

\textbf{23.} 90 करोड़ वर्ष पहले के बाद (खड़िया वलित है) और 2.58 मिलियन वर्ष पहले से पहले (जलोढ़ भ्रंशित नहीं है);
\omterm{def:b1:geological-maps:fold}{वलन} पहले आया, भ्रंशन उसके बाद।

\textbf{24.} भ्रंश F1 (\qty{405}{m}) से \qty{-1000}{m} तक \qty{1405}{m} नीचे जाता है; \ang{60} की \omterm{def:b1:geological-maps:attitude}{नति} पर इसके लिए
नीचे-ढलान की ओर क्षैतिज दूरी $1405/\tan\ang{60} \approx \qty{811}{m}$ चाहिए। कुआँ F1 से लगभग
\textbf{\qty[text-series-to-math=true]{810}{m} पूर्व-दक्षिण-पूर्व} में (दिगंश \ang{120}) खोदा जाना चाहिए।
\end{solution}
